\section{The canonical theta lift and preservation of valuation}
\label{sec:theta-lift-valuation}

We keep the chart and uniformization of
Section~\ref{sec:theta-translate-span}, and write \(F=\overline K\)
and \(\mathcal V=\{c\in F:v(c)\ge0\}\). The two statements used in the
integral spanning argument have different hypotheses. A point killed
by \(\alpha\) has a lift fixing the pulled-back theta function.
An arbitrary finite-order lift preserves every normalized coefficient
valuation bound, even when the theta function being translated is
unrelated to that pullback.

\paragraph{The lift fixing the pullback.}
On the generic abelian variety let \(f=[\alpha]\). The equality
\(f[\overline\alpha]=[\overline\alpha]f=[q]\) implies that \(f\) is
surjective and that its kernel lies in \(A[q]\). Hence \(f\) is an
isogeny. We work over the algebraically closed field \(F\) of
characteristic zero, so this is a separable isogeny.
If \(\mathcal L\) is the ample theta bundle, set
\(\mathcal M=f^*\mathcal L\), with the fixed normalized identification
with the level-\(q\) bundle. For \(x\in\ker f\), \(fT_x=f\).
Pullback of this equality gives
\(T_x^*\mathcal M\simeq\mathcal M\). These identifications satisfy the
group law by functoriality, so they split the theta group above
\(\ker f\). Each fixes \(f^*s\) for every section \(s\), and its order
divides \(q\). This is precisely the pullback construction in Mumford's
level-subgroup correspondence: the isogeny, the descended ample
line bundle, and the identification of its pullback are specified here.
\footnote{\href{https://pazuki.perso.math.cnrs.fr/index_fichiers/Mumford66.pdf}{D.~Mumford,
\emph{On the equations defining abelian varieties I}, \S1,
Proposition~1, pp.~290--291}.}

We describe the compatibility required to use this lift on the
uniformized coefficients. Let \(p\) represent \(x\) in the
uniformization at level \(q\), with lattice coordinate \(y\in X\)
and abelian coordinate \(b\). Since \(f(x)=0\), the torsion
isomorphism gives a period \(z\in X\) with
\(\alpha y+qz=0\). Multiplication by \(\overline\alpha\) and
torsion freeness of \(X\) give \(y=-\overline\alpha z\).
On the abelian part \([\alpha]b=c(-z)\), where \(c\) is the period
map to the abelian part. The generator identifying the CM adjoint
with Rosati is rigidified to one at the origin. Write \(R_\chi\)
for its value on \((c(\chi),b)\). If \(G_{\overline\alpha\chi}\)
is the theta-torsor factor in the chosen representative, then the
multiplication rule for addition of the compensating period \(z\)
gives
\[
 G_{\overline\alpha\chi}R_\chi=\tau(-z,\chi).
\]
Indeed the class of \([\alpha]p\) plus the image of \(z\) is zero.
In the addition formula the cubical factor multiplying
\(\tau(-z,\chi)\tau(z,\chi)\) is the same one in the period
cocycle for \((-z,z)\); cancelling it and using \(\tau(0,\chi)=1\)
gives the displayed identity in common frames.

The coefficient of \(\theta_\alpha\) outside
\(\overline\alpha X\) vanishes. Translation by \(y\) preserves
that subset. At an index \(\overline\alpha\chi\), its translate
has index \(\overline\alpha(\chi-z)\). There are two ways to compare
the corresponding lines: pull back the \((-z)\)-period functional
equation of the original theta function, or translate its pullback
by \(b\). At the origin choose the lift generator so that these
comparisons agree, using the displayed identity and the specified
rigidification of the power-\(q\) bundle. The ratio of the two
cubical identifications at any point is a scalar: the line maps
differ by a global unit on the proper geometrically connected
abelian part. Its value at the origin is one, so they agree
everywhere. This verifies
\(\Phi_x\theta_\alpha=\theta_\alpha\) in the expansion action.
For a representative of a \(q\)-torsion point the \(q\)-th power
of its lift is central, hence acts by a scalar \(c\). As
\(\theta_\alpha\) is primitive and nonzero in our use, this
eigenvalue identity forces \(c=1\). In the version that allows the
zero theta function, any scalar \(q\)-th root of \(c^{-1}\) rescales
the operator to finite order and still fixes zero.
\footnote{\href{https://github.com/ChristianKudeba/alpoge-rh-verification/blob/alpoge-source-7cbf7a0c/src/ReciprocityTheory/ReductiveGroups/ArithmeticGeometry/FirstCohomology/CanonicalLiftUnpacking.lean}{Period representing the zero torsion point};
\href{https://github.com/ChristianKudeba/alpoge-rh-verification/blob/alpoge-source-7cbf7a0c/src/ReciprocityTheory/ReductiveGroups/ArithmeticGeometry/FirstCohomology/CanonicalLiftOrigin.lean}{Rosati normalization and cancellation at the origin};
\href{https://github.com/ChristianKudeba/alpoge-rh-verification/blob/alpoge-source-7cbf7a0c/src/ReciprocityTheory/ReductiveGroups/ArithmeticGeometry/FirstCohomology/CanonicalLiftFrobenius.lean}{comparison at each character};
\href{https://github.com/ChristianKudeba/alpoge-rh-verification/blob/alpoge-source-7cbf7a0c/src/ReciprocityTheory/ReductiveGroups/ArithmeticGeometry/FirstCohomology/CanonLiftPf.lean}{finite-order assertion}.}

\paragraph{The valuation shift for a general lift.}
For a level-\(q\) family \(s\) and \(r\in\mathbb Q\), denote by
\(\mathcal U_r(s)\) the following predicate. For every character
\(\chi\), its coefficient \(s_\chi\) is a finite combination of
base changes of integral sections at \(\chi\), with scalar
coefficients of valuation at least \(r+Q(\chi)/q\).
In particular the earlier \(\mathcal U\) consists of those \(s\)
satisfying \(\mathcal U_0(s)\).

Consider a lift represented by a torsion datum with lattice shift
\(y\), abelian coordinate \(b\), line generator \(w\), and scalar
\(\zeta\in F^\times\); write \(S\) for its operator. Properness
of the abelian part provides an extension \(m\) of \(b\) to
\(\mathcal V\). An integral trivialization of the line for \(w\)
there has ratio \(c_w\in F^\times\) to \(w\), with convention
\(w=c_w\) times the integral generator. We claim
\begin{equation}\label{eq:theta-liftshift}
 \mathcal U_r(s)\quad\Longleftrightarrow\quad
 \mathcal U_{r+C}(Ss),\qquad
       C=v(c_w)+v(\zeta)-Q(y)/q .
\end{equation}

For the index-dependent theta-torsor factor, choose an integral
generator \(G\) of the biextension line at \((m,c(\chi))\); write
the factor as \(t_\chi G\). We compute
\begin{equation}\label{eq:cubical-ratio}
                        qv(t_\chi)=B(y,\chi).
\end{equation}
The torsion relation is \(qm=c(y)\). The period multiplier at
level \(q\) has order \(qQ(y)+B(y,\chi)\): its factors have
orders \(qQ(y)\), \(Q(y+\chi)-Q(y)-Q(\chi)\), and the rigidification
order zero. Denote the scalar of that multiplier in an integral
frame by \(c_0\), the scalar of the \(q\)-th power of the period
generator at \(y\) by \(c_\psi\), and the scalar of the \(q\)-th
power of the rigidification by \(c_1\). Their valuations are
\(qQ(y)+B(y,\chi),qQ(y),0\), respectively.

To see the cancellation, trivialize all factors in a common chart.
Write \(P,R\) for the last two powered generators, \(\tau\) for the
period biextension factor, and \(D\) for the cubical identification
iterated from \(m\) to \(qm\). The family multiplier is
\(\tau P R\), and the theta-torsor chain at \(p\) gives
\((t_\chi G)^q=\tau D\). Write \(P=c_\psi E_\psi\),
\(R=c_1 E_1\), and \(\tau PR=c_0 E\) with the respective integral
generators \(E_\psi,E_1,E\). It follows that
\[
       c_0 A=t_\chi^q c_\psi c_1,\qquad
       A=E E_\psi^{-1} E_1^{-1} D G^{-q}.
\]
All the factors in \(A\) are integral generators of their lines.
Their product is a generator of the trivial line on
\(\operatorname{Spec}\mathcal V\): the tensor factors for \(m\)
and \(qm\) cancel by \(qm=c(y)\) and the cubical identity.
Therefore \(A\in\mathcal V^\times\). Taking valuations in the
previous equality proves \eqref{eq:cubical-ratio}, including its
sign and its factor \(q\).
\footnote{\href{https://github.com/ChristianKudeba/alpoge-rh-verification/blob/alpoge-source-7cbf7a0c/src/ReciprocityTheory/ReductiveGroups/ArithmeticGeometry/DiscreteValuationRings/ExistsValGenHasVal.lean}{Orders of the factors of the period multiplier};
\href{https://github.com/ChristianKudeba/alpoge-rh-verification/blob/alpoge-source-7cbf7a0c/src/Development/AbelianVarietyTheory/Proofs/RingsOfIntegers/ValuationThetaRatioEqBilinear.lean}{cubical trivialization and scalar identity}.}

A coefficient at \(\chi\) is carried to one at \(\chi+y\)
by translation through \(m\) and multiplication by the integral
line isomorphism just described and the scalar \(\zeta c_w t_\chi\).
Translation through \(m\), and its inverse, preserve integral
sections over \(\mathcal V\). They can equivalently be expressed
as combinations of sections over \(T\). Indeed \(T\) is a discrete
valuation ring and \(\mathcal V\subset F\) is torsion free over \(T\),
hence flat. The abelian scheme over \(T\) is proper, thus
quasi-compact and quasi-separated, and each coefficient sheaf is
invertible. Flat base change for \(H^0\) identifies the section
module after extension with the tensor product of the section
module over \(T\) and \(\mathcal V\). This proves the instance of
integral base change needed here, including for a nondiscrete
\(\mathcal V\).
\footnote{\href{https://stacks.math.columbia.edu/tag/0539}{Stacks Project,
Lemma~15.22.10, flatness over a valuation ring};
\href{https://stacks.math.columbia.edu/tag/02KH}{Lemma~30.5.2, flat base change for sections}.}

The gain in the unnormalized lower bound of the coefficient
is \(v(c_w)+v(\zeta)+B(y,\chi)/q\). Subtracting the new
\(Q(\chi+y)/q\) and restoring the old \(Q(\chi)/q\) gives exactly
\(C\). This proves the forward implication in
\eqref{eq:theta-liftshift}; translation by \(-m\) and the inverse
line generator gives the reverse implication. No nonzero
coefficient or attained minimum needs to be chosen. In the usual
analytic notation for a Raynaud extension, the underlying period
functional equation for coefficients is
\(a_{u+\lambda(z)}=t(z,u)c(z)T_{\Phi(z)}^*a_u\), interpreted in
the corresponding line bundles. Foster, Rabinoff, Shokrieh and Soto
give both this formula and the equality of its model norms over
a complete algebraically closed nonarchimedean field. It applies
after passing from \(F\) to its completion in the semiabelian
case here. The coefficient-span calculation above works directly
over \(F\) and supplies \eqref{eq:theta-liftshift} for the chosen
normalization.
\footnote{\href{https://arxiv.org/pdf/1705.08212}{T.~Foster,
J.~Rabinoff, F.~Shokrieh, A.~Soto, \emph{Non-Archimedean and
tropical theta functions}, Proposition~3.19 and Corollary~3.21};
\href{https://github.com/ChristianKudeba/alpoge-rh-verification/blob/alpoge-source-7cbf7a0c/src/ReciprocityTheory/ReductiveGroups/ArithmeticGeometry/DiscreteValuationRings/TWMain.lean}{the full scalar shift in the coefficient filtration}.}

\paragraph{Eliminating the shift for a finite-order lift.}
Suppose \(S^h=1\), \(h\ge1\). For fixed \(s\) the subset
\(D_s=\{r\in\mathbb Q:\mathcal U_r(s)\}\) is downward closed,
and \eqref{eq:theta-liftshift} gives \(D_s=D_s+hC\).
If \(C=0\) the assertion is immediate. If \(C\ne0\), the only
downward-closed subsets of \(\mathbb Q\) with that period are the
empty set and all of \(\mathbb Q\): from one element \(u\),
\(u+j|hC|\) belongs to the set for every \(j\ge0\). Thus in this
case \(D_s+C=D_s\) also. It follows that every finite-order
lift preserves \(\mathcal U_r\) for every \(r\), including positive
depth and zero depth. This proof covers the zero family, whose
valuation bounds are unbounded above.
\footnote{\href{https://github.com/ChristianKudeba/alpoge-rh-verification/blob/alpoge-source-7cbf7a0c/src/ReciprocityTheory/ReductiveGroups/ArithmeticGeometry/DiscreteValuationRings/PeriodicValuation.lean}{Periodicity and exact finite-order invariance}.}

The use of these identities assumes the polarized uniformization,
its torsion isometry and its rigidified cubical period data. Section~\ref{sec:theta-unit-basis} proves the integral rank of the
theta expansion module and the primitivity assertions used by
Section~\ref{sec:theta-translate-span}.
